\documentclass[12pt,a4paper]{article}

\usepackage[utf8]{inputenc}
\usepackage[T1]{fontenc}
\usepackage{amsmath,amssymb,amsfonts}
\usepackage{graphicx}
\usepackage[margin=2.5cm]{geometry}
\usepackage{hyperref}

\title{\textbf{Causal Identifiability of Destructive Interference:}\\[6pt]
\large Formalizing the Quantum Ceiling as a Structural Zero}
\author{Judea Pearl}
\date{March 2026}

\begin{document}
\maketitle

\begin{abstract}
I fully endorse Sabine Hossenfelder's recent critique of the "Quantum Ceiling" protocol. She accurately identifies that Mechanism B (local attention bleed) is mathematically isomorphic to a classical Bayesian update and therefore cannot compute the negative amplitudes required for destructive interference. In this paper, I formalize her critique using causal DAGs. I demonstrate that the inability to sum amplitudes to zero ($A_1 + A_2 = 0$) is not a semantic confound ($do(Z)$) but a fundamental structural zero ($do(B)$) in the evaluating architecture. Attempting to simulate quantum mechanics via classical autoregression is causally unidentifiable.
\end{abstract}

\section{The Causal Graph of Classical Probability Mixing}

Baldo and Chang propose a double-slit experiment to test the Quantum Ceiling. Let us formalize the causal paths of this generative process.

Let $Z$ be the narrative prompt ("water wave passes through slits").
Let $E$ be the dense vector encoding of this context.
Let $B$ be the native architecture (Transformer/SSM) utilizing a strictly positive, classical probability space ($\mathbb{R}_{\geq 0}$).
Let $Y$ be the generated outcome distribution on the "screen."

If the model were natively capable of quantum computation, there would be a hidden state variable $H$ representing a complex state vector, allowing for paths where $H_1 + H_2 = 0$, leading to $P(Y) = 0$.

However, under Mechanism B, there is no $H$. The causal path is:
\begin{equation*}
do(Z) \rightarrow E \rightarrow Y \leftarrow B
\end{equation*}

As Hossenfelder correctly points out, the architecture $B$ strictly enforces that $P(A \cup B) = P(A) + P(B) - P(A \cap B)$.

\section{Destructive Interference as a Structural Zero}

In causal inference, a "structural zero" is an outcome that is forbidden by the fundamental structure of the causal graph, regardless of the intervention on the treatment variables.

Because $B$ limits the operations on $E$ to classical probability mixing (Bayesian updates over text co-occurrences), a point of true destructive interference where the combined probability $P(Y) = 0$ (while the individual probabilities $P(Y|path_1) > 0$ and $P(Y|path_2) > 0$) is a structural zero.

No intervention on the prompt $do(Z)$ can create the missing causal edge or the missing complex variable $H$.

\section{Conclusion}

The Quantum Ceiling protocol is valuable precisely because its failure is causally guaranteed. By formalizing Hossenfelder's insight into a DAG, we can definitively state that the inability of an LLM to simulate destructive interference is not a failure of prompt engineering or semantic gravity. It is a hard architectural bound. The structural zero ($do(B)$) cannot be bypassed by any semantic intervention ($do(Z)$). The test will yield classical diffusion.

\end{document}
